% concatenated from Wang-Wang-Rasila_2026__3.21_.tex
%------------------------------------------------------------------------------
% Here please write the date of submission of paper or its revisions:
%------------------------------------------------------------------------------
%
\documentclass[11pt, reqno]{amsart}
%%%%%%%Begins%%%%Use the first line with .eps. Use the second one when you use .pdf for pictures
%\usepackage{graphicx}
%\usepackage{lineno,hyperref}
%%%%%%Ends%%%%%%%%

\usepackage{ifpdf}
\ifpdf
  \usepackage[hyperindex,pagebackref]{hyperref}%
\else
  \expandafter\ifx\csname dvipdfm\endcsname\relax
    \usepackage[hypertex,hyperindex,pagebackref]{hyperref}
  \else
    \usepackage[dvipdfm,hyperindex,pagebackref]{hyperref}
  \fi
\fi

\usepackage{cases}
%\usepackage{txfonts}
%\usepackage{amsmath}
\usepackage[square, comma, sort&compress, numbers]{natbib}%参考文献连续编号
\usepackage{float}
\usepackage{pifont}
\usepackage{bbding}
\usepackage{amssymb}
\usepackage{mathrsfs}
\usepackage{subfigure}
\usepackage{caption}
\usepackage{geometry}
\usepackage{amsmath, amsthm, amscd, amsfonts, amssymb, graphicx, color}
\usepackage{lineno}

\usepackage{ifpdf}



\textheight 21.0truecm \textwidth 13.7truecm
\setlength{\oddsidemargin}{0.35in}\setlength{\evensidemargin}{0.35in}
\bibliographystyle{alpha}
\setlength{\topmargin}{-.5cm}

\newtheorem{theorem}{Theorem}
\newtheorem{thm}{Theorem}
\newtheorem{lemma}{Lemma}
\newtheorem{prop}{Proposition}
\newtheorem{prob}{Problem}
\newtheorem{preim}{Lemma}
\newtheorem{cor}{Corollary}
\newtheorem{con}{Conjecture}
\newtheorem{dfn}{Definition}
\newtheorem{rem}{Remark}
\newtheorem{example}{Example}
%\renewcommand{\thethm}{\Roman{thm}}
\renewcommand{\thethm}{\Alph{thm}}
\numberwithin{equation}{section}
\newcommand{\nbd}{\nobreakdash}





\newcommand{\abs}[1]{\left\vert#1\right\vert}
\newcommand{\U}{\mbox{$\mathbb{U}$}}
\newcommand{\R}{\mbox{$\mathbb{R}$}}
\newcommand{\N}{\mbox{$\mathbb{N}$}}
\newcommand{\Q}{\mbox{$\mathbb{Q}$}}
\newcommand{\C}{\mbox{$\mathbb{C}$}}
\newcommand{\Z}{\mbox{$\mathbb{Z}$}}
\newcommand{\K}{\mbox{$\mathbb{K}$}}
\newcommand{\I}{\mbox{$\mathbb{I}$}}
\newcommand{\D}{\mbox{$\mathbb{D}$}}
\newcommand{\ps}{\pi/\sin\frac{\pi}{p}}
\newcommand{\nsum}{\sum^\infty_{n=1}}
\newcommand{\msum}{\sum^\infty_{m=1}}
\newcommand{\mmsum}{\sum^{k-1}_{m=1}}
\newcommand{\fint}{\int^\infty_0f(t)\,dt}
\newcommand{\fkint}{\int^\infty_kf(t)\,dt}
\newcommand{\rlam}{\lambda(r)}
\newcommand{\ffm}{\frac{1}{1+m}\left(\frac{1}{m}\right)^{\frac{1}{r}}}
\DeclareMathOperator{\td}{d\mspace{-2mu}}
%\renewcommand{\baselinestretch}{1.5}

\makeatletter
\@namedef{subjclassname@2020}{\textup{2020} Mathematics Subject Classification}
\makeatother

\date{\today}
\begin{document}
\setcounter{page}{1}




\title[Harmonic $K$-quasiconformal Koebe functions]
{Harmonic $K$-quasiconformal Koebe functions: construction and application to Pavlovi\'c's problem
%and \\ \vskip.05in harmonic Hardy spaces
}

\author[Zhi-Gang Wang, Xiao-Yuan Wang, Antti Rasila, and Jia-Le Qiu
]{Zhi-Gang Wang, Xiao-Yuan Wang, Antti Rasila, and Jia-Le Qiu
}

\vskip.20in
\address{\noindent Zhi-Gang Wang\vskip.05in
School of Mathematics and Statistics, Hunan
	First Normal University, Changsha 410205, Hunan, P. R. China.}

%\address{\noindent 
		%School of Mathematics and Statistics, Changsha University of Science and Technology,
		%Changsha 410114, Hunan, P. R. China.}
 
\vskip.05in
\email{\textcolor[rgb]{0.00,0.00,0.84}{wangmath$@$163.com}}






\vskip.05in
\address{\noindent Xiao-Yuan Wang \vskip.05in
	School of Mathematical Sciences, Liaocheng University, Liaocheng
	252059, Shandong, P. R. 
	China.}

%\address{\noindent Graduate School of Information Sciences, Tohoku University, Sendai 980-8579, Japan.}
\vskip.05in
\email{\textcolor[rgb]{0.00,0.00,0.84}{mewangxiaoyuan$@$163.com}}
 
\vskip.05in
\address{\noindent Antti Rasila \vskip.05in
	Department of Mathematics with Computer Science, Guangdong Technion-Israel Institute
of Technology, 241 Daxue Road, Shantou 515063, Guangdong, P. R.
	China.}

\address{\noindent
Department of Mathematics, Technion-Israel Institute of Technology, Haifa 3200003, Israel.}
\vskip.05in
\email{\textcolor[rgb]{0.00,0.00,0.84}{antti.rasila$@$iki.fi}; \textcolor[rgb]{0.00,0.00,0.84}{antti.rasila$@$gtiit.edu.cn}}


\vskip.05in
	\address{\noindent Jia-Le Qiu \vskip.03in
		School of Mathematics and Statistics, Changsha University of Science and Technology,
		Changsha 410114, Hunan, P. R. China.}
  \vskip.05in
	\email{\textcolor[rgb]{0.00,0.00,0.84}{qiujiale2023$@$163.com}}




%\vskip.05in
%\email{\textcolor[rgb]{0.00,0.00,0.84}{wangmath$@$163.com}}


%\address{\noindent Zhi-Hong Liu\vskip.05in
%College of Science, Guilin University of Technology, Guilin 541004, Guangxi, P. R. China.}
%\email{\textcolor[rgb]{0.00,0.00,0.84}{liuzhihongmath$@$163.com}}


\subjclass[2020]{31A05, 30C55, 30C62, 30H10.}

\keywords{Harmonic $K$-quasiconformal Koebe function; harmonic $K$-quasiconformal mapping; harmonic Hardy space; harmonic Schwarzian derivative.}



\begin{abstract} 

We first construct the harmonic $K$-quasiconformal Koebe functions, filling a long-standing foundational gap in geometric function theory. This construction provides a unified parametric candidate extremal function framework for conformal mappings, quasiconformal mappings, and harmonic mappings, and we formulate related conjectures for the extremal theory of harmonic $K$-quasiconformal mappings. By combining this construction with Astala and Koskela's $H^p$-theory for quasiconformal mappings, we establish a sharp result concerning the optimal order of harmonic $K$-quasiconformal mappings with bounded Schwarzian norm in harmonic Hardy spaces. Motivated by the work of Chuaqui, Hernández, and Mart\'in’s [Math. Ann. 367, 1099–1122 (2017)], this result gives a partial solution to Pavlovi\'c's  2014 open problem on the embeddings of harmonic quasiconformal mappings into Hardy spaces, and outlines a path toward its complete solution.
\end{abstract} \maketitle



%\tableofcontents

\section{Introduction and statement of the main result}\label{s1} 

Following the seminal work of Clunie and Sheil-Small \cite{cs}, univalent harmonic mappings have emerged as a central topic at the interface of complex analysis, differential geometry, geometric topology, Teichm\"{u}ller theory, nonlinear dynamics, and numerical analysis. Despite remarkable progress over recent decades, several fundamental problems remain strikingly open. Most notably, sharp analogues of Bieberbach's conjecture (de Branges's theorem \cite{de}) and the Koebe one-quarter covering theorem, cornerstones of classical conformal mapping theory, still lack definitive solutions in the harmonic quasiconformal setting. Within geometric function theory, a long-standing bottleneck persists: there exists no effective framework for constructing and rigorously establishing Koebe-type conjectures for harmonic $K$-quasiconformal (\textbf{HQC}) mappings. This critical gap has severely hampered the systematic development of extremal theory and canonical example construction in the field. 

Linear properties, the maximum modulus principle, qualitative results in quasiconformal geometry and the like extend trivially to the setting of harmonic quasiconformal geometric function theory; however, the essence of classical geometric function theory, centered on extremal functions, coefficient conjectures, sharp distortion inequalities and Loewner theory, cannot be extended trivially at all, and a completely new theoretical framework must be established. This is precisely the core value and research significance of harmonic quasiconformal geometric function theory.

%Chuaqui, Hern\'{a}ndez, and Mart\'{i}n \cite{chm} studied the affine and linearly invariant families of planar harmonic mappings in the unit disk. Liu and Ponnusamy \cite{lp}, Ponnusamy, Qiao, and Wang \cite{pqw} have derived various properties of uniformly locally univalent harmonic mappings. Wang, Liu, Rasila, and Sun \cite{wlrs} constructed a counterexample to make a negative answer to a question posed by Bharanedhar and Ponnusamy \cite{bp}. Arbel\'{a}ez, Hern\'{a}ndez, and Sierra \cite{ahs} discussed Lower and upper order of harmonic mappings. We also note that Schoen and Yau \cite{sy} have studied univalent harmonic maps between surfaces.
%S\`{e}te and Zur \cite{sz} designed a Newton method to process the numerical problems in planar harmonic mappings.



%For $a\in\C$ and $r>0$, let $\D(a,r):=\{z:\,\abs{z-a}<r\}$. Particularly, let $\D_r$ denote
%the disk $\D(0,r)$ and $\D$ for the open unit disk $\D_1$.

Let $\mathcal{H}$ denote the class of complex-valued harmonic functions
$f=h+\overline{g}$ in the open unit disk $\D$, normalized by the conditions $$f(0)=f_{z}(0)-1=0,$$ which have the
form
\begin{equation}\label{111}
f(z)=z+\sum_{n=2}^{\infty}a_nz^n+\overline{\sum_{n=1}^{\infty}b_nz^n}.
\end{equation}

Denote by $\mathcal{A}$ the class of normalized analytic functions, which is a subclass of
$\mathcal{H}$ with $g(z)\equiv 0$.
Let $\mathcal{S_{\mathcal{H}}}$ be the class of harmonic functions $f\in\mathcal{H}$
that are univalent and sense-preserving in $\D$. Moreover, we denote by $\mathcal{S}^{0}_{\mathcal{H}}$ the subclass of $\mathcal{S_{\mathcal{H}}}$ with the additional condition $f_{\overline{z}}(0)=0$. Since there exist reciprocal transformations between the classes $\mathcal{S_{\mathcal{H}}}$
and $\mathcal{S}^{0}_{\mathcal{H}}$, we usually focus our attention on the class $\mathcal{S}^{0}_{\mathcal{H}}$.  Clearly, the classical class $\mathcal{S}$ of normalized univalent analytic functions is a subclass of the class
$\mathcal{S}_\mathcal{H}^0$. %with $g(z)\equiv 0$.




\subsection{Harmonic quasiconformal mappings}
We say that $f$ belongs to the class $\mathcal{S}_\mathcal{H}(K)$ of \textbf{HQC} mappings, where $K\ge 1$ is a constant, if $f\in\mathcal{S}_\mathcal{H}$ and its dilatation defined by $$\omega:=\frac{g'}{h'}$$ satisfies the condition 
$\abs{\omega}\leq k$, where $k\in [0,1)$ is given by
$$
k:=\frac{K-1}{K+1}\quad(K\geq 1).
$$ A function $f$ is called a harmonic quasiconformal mapping, if it belongs to $\mathcal{S}_\mathcal{H}(K)$ for some $K\ge 1$.

For convenience, we define $$\mathcal{S}^0_\mathcal{H}(K):=\mathcal{S}_\mathcal{H}(K)\cap\mathcal{S}^0_\mathcal{H}.$$
Note that $$\mathcal{S}^0_\mathcal{H}(1)=\mathcal{S}_\mathcal{H}(1)=\mathcal{S}.$$


%The class of harmonic quasiconformal mappings was first studied by Martio in \cite{m}, and his topic has attracted  substantial interest since then.
%For recent results on harmonic quasiconformal mappings, we refer the reader to the following works.
%Chen and Ponnusamy \cite{cp2}, \cite{cp3}, \cite{cp}, Arbel\'{a}ez,  Chuaqui and Sierra \cite{acs} studied John disk, radial length, radial John disk and Koebe type theorems associated with harmonic quasiconformal
%mappings. 
%Chen and Fang \cite{cf} derived a Schwarz-Pick inequality for harmonic quasiconformal
%mappings.
%Huang \cite{h} discussed the characterizations of harmonic quasiconformal mappings on the upper half-plane.
%Fu and Huang \cite{fh} considered the family of harmonic $K$-quasiconformal mappings from unit disk onto half
%planes.  
In recent years,
Kalaj \cite{ka} obtained the extended form of Lindel\"{o}f theorem for harmonic quasiconformal mappings.
Kalaj \cite{ka2} established sharp Riesz-type inequalities for harmonic mappings on the open unit disk in the complex plane.
Wang, Shi, and Jiang \cite{wsj} investigated harmonic quasiconformal mappings associated with asymmetric vertical strips. 
Sun, Rasila, and Jiang \cite{srj} considered linear combinations of harmonic quasiconformal mappings convex
in one direction.
Chuaqui, Hern\'{a}ndez, and Mart\'{i}n \cite{chm}
studied the affine and linear invariant families of
harmonic mappings, and obtained the sharp order of harmonic mappings with bounded Schwarzian norm.
Partyka, Sakan, and Zhu \cite{psz}  provided various properties of harmonic quasiconformal mappings with the
convex holomorphic part.
%Gjokaj and Kalaj \cite{gk} derived the properties of harmonic quasiconformal mappings between the unit ball
%and a spatial domain. 
Liu and Zhu \cite{lz} studied the Riesz conjugate functions theorem for harmonic quasiconformal
mappings. %Liu and Ponnusamy \cite{lp2} obtained Bohr radius for harmonic quasiconformal
%mappings. 
  %Wang, Huang, Liu and Kargar \cite{whlk} obtained several properties of close-to-convex harmonic quasiconformal mappings whose
%analytic parts are starlike mappings. 
Todor\v{c}evi\'{c}'s monograph \cite{t} gives recent developments on higher-dimensional harmonic quasiconformal mappings and hyperbolic type metrics.

%Zhu \cite{z} derived some coefficient estimates for harmonic $K$-quasiconformal mappings. 
%Ha and Kim \cite{hk} characterized the %Gaussian curvature estimates over maximal %graphs of harmonic $K$-quasiconformal mappings.


%Here, we present the following %intuitive relationships among the classes $\mathcal{S}$, ${\mathcal{S}_\mathcal{H}^0}$, $K$-$QC$ and ${\mathcal{S}_\mathcal{H}^0(K)}$.
%where $K$-$QC$ denotes the familiar class of $K$-quasiconformal mappings.
%{\begin{figure}[ht]
%\centering
%\includegraphics[width=2.91in]{Figure1.pdf}
%\caption{\small Relationships among the classes $\mathcal{S}$, ${\mathcal{S}_\mathcal{H}^0}$, $K$-$QC$ and ${\mathcal{S}_\mathcal{H}^0(K)}$.
%\begin{center}
%${S_H^0(1)}=:S$ and $\lim\limits_{K\rightarrow +\infty}{S_H^0(K)}=:{S_H^0}$.\end{center}
%}\label{fig1}
%\end{figure}}

\subsection{Hardy spaces for harmonic and quasiconformal mappings}
Let $0<p \leq \infty$. A complex-valued function $f$ in $\D$ is said to be in the {\it Hardy space} for $p$, if the integral mean
 \[M_p(r,f):= \begin{cases}
\left(\frac{1}{2\pi}\int_0^{2\pi}\abs{f(re^{i\theta})}^pd\theta\right)^{\frac{1}{p}}&(0<p<\infty),\\ \ \
\sup\limits_{{\abs{z}=r}}\abs{f(z)}&(p=\infty),
\end{cases}\]
is defined for all $r\in(0,1)$, and
$$\sup\limits_{0<r<1}M_p(r, f)<\infty.$$ 

We record the following special cases of this definition. The classical Hardy space of analytic functions is denoted by $H^p$ 
%(for details, see \cite{du70,d,ka2})
and the one of $K$-quasiconformal mappings \cite{ak} is denoted by $H^p_K$. Similarly, we say that a harmonic function $f$ is in a \textit{harmonic Hardy space} $h^p$, if it is harmonic and satisfies the above condition. We denote the space of \textbf{HQC} mappings by $h^p_K$. 

\subsection{The order of a Hardy space} %involving $K$-quasiconformal mappings and univalent harmonic mappings

For a class $\mathcal{G}$ of complex-valued functions in $\D$, we define its {\it order} as a constant $\gamma\geq 0$ such that $\mathcal{G}$ is contained the corresponding Hardy space $H^p$ for all $p\in(0,\gamma)$.

We recall the following special case of the result due to Astala and Koskela \cite[Theorem 3.2]{ak}.% (see also \cite{bkl,bk}).

\vskip.10in

\begin{thm} \label{thmB}
 If $f$ is a $K$-quasiconformal mapping in $\D$, then $f\in H^p_K$ for $$0<p<\frac{1}{2K}\quad(K\geq1),$$ and the constant $1/(2K)$ is sharp.
 \end{thm}

\vskip.10in

In 1990, Abu-Muhanna and Lyzzaik \cite[Theorem 2]{al} established the following result showing that for a harmonic mapping $f=h+\overline{g}$, the functions $h$ and $g$ belong to the Hardy space $H^p$, and $f$ belongs to harmonic Hardy spaces $h^p$, respectively, for a suitable choice of $p$.

\vskip.10in

\begin{thm} \label{thmC}
 If $f=h+\overline{g}\in \mathcal{S}_\mathcal{H}$, then $h,\, g\in H^p$ and $f\in h^p$ for every $p$ with $$0<p<\frac{1}{(2\beta+2)^{2}},$$
where $$\beta:=\sup\limits_{f\in\mathcal{S}_\mathcal{H}}\abs{a_2}.$$
\end{thm}


In 1996, the order $$\frac{1}{(2\beta+2)^{2}}$$ in Theorem \ref{thmC} was improved to $1/\beta^{2}$ by Nowak in \cite[Theorem 1]{n}, where she conjectured that the order is actually $1/\beta$. Moreover, she obtained the sharp results $$f\in h^p\quad\left(0<p<\frac{1}{2}\right)$$ for the class $\mathcal{K}_\mathcal{H}$ of convex harmonic mappings, and 
$$f\in h^p\quad\left(0<p<\frac{1}{3}\right)$$ for the class $\mathcal{C}_\mathcal{H}$ of close-to-convex harmonic mappings. Even today, proving  Nowak's conjecture is seen as a 
\textit{challenging} problem.



We also note that Aleman and Mart\'{\i}n \cite{am} proved that convex harmonic mappings are not necessarily in $h^{1/2}$, which provided a negative answer
to a question raised by Duren \cite[p. 153, Sect. 8.5]{d}. Laitila, Nieminen, Saksman, and Tylli \cite{ln} considered the rigidity of composition operators involving harmonic Hardy spaces $h^p$. Melentijevi\'{c} and Bo\v{z}in \cite{mb} obtained sharp Riesz-Fej\'{e}r inequality for harmonic Hardy spaces $h^p$.

An analytic function $h$ in $\D$ has valency $n$ if $h$ takes no value
more than $n$ times. More generally, for a function $h$ analytic in $\D$, let $W(R)$ denote
the area (regions covered multiply being counted multiply) of the image of $\D$ under
$h$ that lies in the closed disk $\abs{w}\leq R$. If $W(R)\leq n\pi R^2$ for all $R>0$, where $n$ is a positive number, then we say that $h$ has mean valency $n$ (cf. \cite{ds}).


In 2024, Das and Sairam Kaliraj \cite[Theorem 1]{ds}
proved that Nowak's conjecture is true with the additional condition that $h^\prime$ has a \textit{finite mean valency}. %of analytic part of $f$. 
Noting that the conjectured $\beta$ of the class $\mathcal{S}_\mathcal{H}$ is equal to $3$ (see \cite{cs}), and it is consequently conjectured by Duren \cite[p. 152, Sect. 8.5]{d} that $$\mathcal{S}_\mathcal{H}\subset h^{{1}/{3}}.$$%\ {\rm and}\ \mathcal{S}_\mathcal{H}^0\subset h^{{2}/{5}}.$$

\subsection{Affine and linear invariant families}
For $f=h+\overline{g}\in\mathcal{S}_\mathcal{H}$, following the terminology of \cite{s}, a family is said to be an \textit{affine
and linear invariant family} ({AL} family) if it is closed under both the \textit{Koebe transform}
$$
K_{\zeta}(f)(z):=\frac{f\left((z+\zeta)/\left(1+\overline{\zeta} z\right)\right)-f(\zeta)}{\left(1-\abs{\zeta}^2\right)h'(\zeta)}
\quad(\abs{\zeta}<1)
$$
and \textit{affine changes}
$$A_{\varepsilon}(f)(z):=\frac{f(z)-\overline{\varepsilon f(z)}}{1-\overline{\varepsilon}g'(0)}
\quad(\abs{\varepsilon}<1).$$

In 1990, Sheil-Small \cite{s} %(see also \cite{ahs}) 
offered an in-depth study of affine and linear invariant families $\mathcal{F}$ of
harmonic mappings $f=h+\overline{g}$ in $\D$. The {\it order} of an AL family is given by
$$\alpha(\mathcal{F}):=\sup\limits_
{f\in\mathcal{F}}
\abs{a_2(f)}=\frac{1}{2}
\sup\limits_
{f\in\mathcal{F}}
\abs{h''(0)}.$$
Indeed, many important characterizations of the AL family $\mathcal{F}$ are determined
by its {\it order}. Clearly, the class $\mathcal{S_{\mathcal{H}}}$ is a typical example of the affine and linear invariant families.

\subsection{Harmonic pre-Schwarzian and Schwarzian derivatives}
The classical pre-Schwarzian and Schwarzian derivatives of locally univalent analytic functions $h$ (cf. \cite{du}),
which are, respectively, defined by
$$P_h(z):=\frac{h''(z)}{h'(z)}$$
and
$$S_h(z):=\left(\frac{h''(z)}{h'(z)}\right)'-\frac{1}{2}\left(\frac{h''(z)}{h'(z)}\right)^2.$$

%We note that Ali and Pal \cite{ap}, and Wang, Li and Fan \cite{wlf} discussed the Schwarzian norm estimates for Janowski convex functions.
%Liao and Wu \cite{lw} considered the exact meromorphic solutions of Schwarzian differential equations.

For a locally univalent sense-preserving harmonic mapping $f=h+\overline{g}$ in $\D$, Hern\'{a}ndez and Mart\'{i}n \cite{hm} introduced
the pre-Schwarzian and Schwarzian derivatives of $f$ as follows:
$$P_f:=\left(\log J_f\right)_z=P_h-\frac{\omega'\,\overline{\omega}}{1-\abs{\omega}^2}$$ and
\begin{align*}\begin{split}S_f&:=\left(\log J_f\right)_{zz}-\frac{1}{2}\left[\left(\log J_f\right)_z\right]^2\\&\, \, =S_h+\frac{\overline{\omega}}{1-\abs{\omega}^2}\left(\frac{h''}{h'}\omega'-\omega''\right)
-\frac{3}{2}\left(\frac{\omega'\,\overline{\omega}}{1-\abs{\omega}^2}\right)^2,\end{split}\end{align*}
where $$J_f:=\abs{f_z}^2-\abs{f_{\overline{z}}}^2$$ is the Jacobian of $f$.
The corresponding pre-Schwarzian norm $\|P_f\|$ and Schwarzian norm
$\|S_f\|$ of $f$ are defined, respectively, by
$$\|P_f\|:=\sup\limits_
{z\in\mathbb{D}}
\abs{P_f(z)}\left(1-\abs{z}^2\right)$$
and
$$\|S_f\|:=\sup\limits_
{z\in\mathbb{D}}
\abs{S_f(z)}\left(1-\abs{z}^2\right)^2.$$

We note that Chuaqui, Duren, and Osgood \cite{cdo} have given an alternative definition of harmonic Schwarzian derivative
with restriction on its dilatation.  
%Efraimidis, Hern\'{a}ndez and Mart\'{i}n \cite{ehm} derived Ahlfors-Weill extensions for harmonic mappings involving harmonic Schwarzian derivative.
%Hern\'{a}ndez and Mart\'{i}n \cite{hm15} confirmed the univalency of locally univalent harmonic mapping with bounded harmonic Schwarzian norm.
%Graf \cite{g} has obtained several estimates on the harmonic pre-Schwarzian and Schwarzian norms of locally harmonic mappings.
%Hern\'{a}ndez and Mart\'{i}n \cite{hm22} characterized the harmonic M\"{o}bius transformations associated with harmonic Schwarzian norm.
%Kanas, Maharana and Prajapat \cite{kmp} studied the norm of the pre-Schwarzian derivative of harmonic mappings.
%Hu and Fan \cite{hf} obtained a criterion for %univalency and quasiconformal extension for %harmonic
%Techm\"{u}ller mappings by using harmonic %Schwarzian derivative. 
%Efraimidis \cite{e} has derived several criteria for univalence and quasiconformal extension for harmonic mappings on planar domains by using 
%harmonic Schwarzian derivative. 


\subsection{A problem of Pavlovi\'{c}}
Because univalent harmonic mappings are not necessarily \textit{quasiconformal}, the constant $1/(2K)$ in Theorem \ref{thmB} may not be sharp for the class of harmonic $K$-quasiconformal mappings. Based on Theorems \ref{thmB} and \ref{thmC}, Pavlovi\'{c} \cite[p. 315, Problem 10.2]{pav} in 2014 proposed the following natural problem.


\begin{prob}\label{p1}
{Find the sharp $$p_0=p_0(K)$$ such that every harmonic $K$-quasiconformal mapping belongs to $h^p_K$ for $p<p_0$.}
\end{prob}

\begin{rem}\label{r02}
{\rm One of keys to solve the above problem %Problem \ref{p1} 
is determination of $\sup\limits_{f\in\mathcal{S}_\mathcal{H}(K)}\abs{a_2}$ of the family $\mathcal{S}_\mathcal{H}(K)$.
But, at present, it seems that this is a \textit{hard} problem (see \cite[Proposition 1.6]{cp} and 
Conjecture \ref{conj1} in Sect. \ref{s3}%and Remark \ref{r01} 
).}
\end{rem}

Curiously, Das, Huang, and Rasila \cite[Theorem 3]{dhr} established the optimal orders for convex and close-to-convex harmonic $K$-quasiconformal mappings to belong to harmonic Hardy spaces, respectively; notably, these results are independent of the quasiconformal constant $K$, i.e.,
$$f\in h^p\quad\left(0<p<1\right)$$ for the class $\mathcal{K}_\mathcal{H}(K)$ of 
convex harmonic quasiconformal mappings, and 
$$f\in h^p\quad\left(0<p<\frac{1}{2}\right)$$ for the class $\mathcal{C}_\mathcal{H}(K)$ of 
close-to-convex harmonic quasiconformal mappings. %(or $\mathcal{S}^*_\mathcal{H}(K)$).

%but determination of the order of the family $\mathcal{S}^0_\mathcal{H}(K)$ is also a challengeable task (see \cite{cp}).
Note that Ponnusamy, Qiao, and Wang \cite[Theorem 4.2]{pqw} also obtained a sharp result regarding uniformly locally univalent harmonic quasiconformal mappings belong to harmonic Hardy space.

 Denote by  $\widehat{\mathcal{S}_\mathcal{H}}(K)$ the subclass of $\mathcal{S}_\mathcal{H}(K)$, %whose members %preserve \textit{affine and linear invariant}, and 
where $h'$ has a \textit{finite mean valency}. We now give the following characteristic of the class $\widehat{\mathcal{S}_\mathcal{H}}(K)$, which will play a crucial role in our study.

\begin{prop}\label{prop1}
\textit{Suppose that $f=h+\overline{g}\in\widehat{\mathcal{S}_\mathcal{H}}(K)$. If $$\phi(K):=\sup\limits_{f\in\widehat{\mathcal{S}_\mathcal{H}}(K)}\abs{a_2}.$$
Then \begin{enumerate}
\item for $\phi(K)\leq 2K$, we have $$f\in h^p\quad \left(0<p<\frac{1}{2K}\right);$$
\vskip.05in
\item for $\phi(K)>2K$, we have $$f\in h^p\quad \left(0<p<\frac{1}{\phi(K)}\right).$$
\end{enumerate}
Both of these bounds are sharp.}
\end{prop}

%Observe that $\varphi(K)\geq 1$ is a continuous function of $K$ in $[1,+\infty)$,
\begin{proof}
By the fact that $\phi(K)$ is a real function with respect to $K$,
$\phi(K)\geq1$ for $K\geq1$ (see \cite{p}), using the similar method as in  Das and Sairam Kaliraj \cite[Theorem 1]{ds}, 
we confirm that if $f\in\widehat{\mathcal{S}_\mathcal{H}}(K)$, %preserve affine and linear invariant, 
then $f\in h^p$ for $$0<p<\frac{1}{\phi(K)}.$$
Note that, by Theorem \ref{thmB} and taking the value of $$\min\limits_{K\geq 1}\left\{\frac{1}{\phi(K)},\ \frac{1}{2K}\right\}.$$ We obtain the assertion of
Proposition \ref{prop1}.
\end{proof}

For convenience, we denote the subclass of ${\mathcal{S}_\mathcal{H}}(K)$ (resp. $\widehat{\mathcal{S}_\mathcal{H}}(K)$) with
\begin{equation*}
 \|S_f\|\leq \lambda%\quad (\lambda\geq 0)
\end{equation*}
by ${\mathcal{S}_\mathcal{H}}(K,\lambda)$ (resp. $\widehat{\mathcal{S}_\mathcal{H}}(K,\lambda)$), where $K\geq1$ and $\lambda\geq0$. %and $$\widehat{\mathcal{S}^0_\mathcal{H}}(K,\lambda):=\widehat{\mathcal{S}_\mathcal{H}}(K,\lambda)\cap\mathcal{S}^0_\mathcal{H}.$$
Now, we establish the following relationship between
the family $\widehat{\mathcal{S}_\mathcal{H}}(K,\lambda)$ of \textbf{HQC} mappings with bounded Schwarzian norm and harmonic Hardy spaces $h^p$.


\begin{theorem}\label{t1}
Let $f\in\widehat{\mathcal{S}_\mathcal{H}}(K,\lambda)$ with $K\geq1$ and $\lambda\geq0$. %such that $f$ preserves affine invariance. 
Then
\begin{enumerate}
\item for $0\leq\lambda\leq6$ and $K\geq1$, we have $$f\in h^p\quad \left(0<p<\frac{1}{2K}\right);$$
\vskip.05in
\item for $\lambda>6$ and $K\geq K_1>1$, we have $$f\in h^p\quad \left(0<p<\frac{1}{2K}\right);$$
\vskip.05in
\item for $\lambda>6$ and $1\leq K<K_1$, we have $$f\in h^p\quad \left(0<p<\frac{1}{\varphi(K,\lambda)}\right),$$
\end{enumerate}
where
\begin{equation}\label{1}\varphi(K,\lambda):=\sup\limits_{f\in\widehat{\mathcal{S}_\mathcal{H}}(K,\lambda)}\abs{a_2}
=\sqrt{1+\frac{\lambda}{2}+\frac{1}{2}\left(\frac{K-1}{K+1}\right)^2}+\frac{K-1}{2\left({K+1}\right)},\end{equation}
and 
$K_1$ is the unique solution in $(1,+\infty)$ of the equation
\begin{equation}\label{2}16K^4+24K^3-(2\lambda-11)K^2-2(2\lambda-1)K-2\lambda-5=0.\end{equation}
All of these results are sharp.
\end{theorem}


\begin{rem}
{\rm For $K=1$ and $0\leq\lambda\leq6$, 
 the first part of Theorem \ref{t1} matches the classical analytic case, as originally examined by Pommerenke \cite{p} in 1964.
%the first part of Theorem \ref{t1} coincides with the classical analytic case studied by Pommerenke \cite{p} %in 1964.
}
\end{rem}

%\begin{rem}
%{\rm Indeed, we consider that the parameter $\lambda$ in Theorem \ref{t1} should be restricted to the condition $\lambda\in[0,19/2)$ (see Remark \ref{r3}). Furthermore, we conjecture that the condition that $h'$ has a finite mean valency in Proposition \ref{prop1} and Theorem \ref{t1} can be removed. This is closely related to Nowak's conjecture (cf. \cite{n}). }
%\end{rem}


This paper employs the weighted shearing technique to construct the first harmonic $K$-quasiconformal Koebe functions for the class, resolving a long-standing core open problem in the field. The constructed harmonic quasiconformal Koebe-type function serves as a candidate extremal function for this class, and a rigorous proof for the corresponding coefficient conjecture remains elusive to date, its difficulty being analogous to that of the classical Bieberbach conjecture. This construction enables the systematic extension of classical results from conformal mapping theory to the harmonic quasiconformal setting, thereby inaugurating a new systematic research program in geometric function theory.
We also establish the sharp order for harmonic $K$-quasiconformal mappings with bounded Schwarzian norm. Furthermore, employing a fundamental result on the $H^p$ theory for quasiconformal mappings due to Astala and Koskela \cite{ak}, we determine the optimal order for the family of harmonic $K$-quasiconformal mappings with bounded Schwarzian norm that lie in harmonic Hardy spaces. This partially resolves an open problem posed by Pavlovi'{c} \cite{pav} in 2014.




\vskip.20in

\section{Preliminaries}

In this section, we provide several required preliminary results.

\subsection{Estimation of $\abs{a_2}$ of a univalent harmonic mapping}
The problem of determining the sharp upper bound of $|{a_2}|$ in the class
$\mathcal{S}^0_\mathcal{H}$
is still an \textit{open} problem. Clunie and Sheil-Small \cite{cs} once proved that 
$$\abs{a_2}<12172.$$
Later, Sheil-Small \cite{s} improved it to $$\abs{a_2}<57.$$ The estimate $$\abs{a_2}<49$$ was
proved by Duren \cite[p. 96, Sect. 6.3]{d}. 
To our knowledge, the latest bound is $$\abs{a_2}<20.9197,$$ which was established
by Abu-Muhanna, Ali, and Ponnusamy \cite{aap} in 2019. 
It has been conjectured by Clunie and Sheil-Small \cite{cs}  that $$\abs{a_2}\leq\frac{5}{2},$$
where the equality is reached by the \textit{harmonic Koebe function} $\mathbb{K}$ defined by
$$\mathbb{K}(z):=\frac{z-\frac{1}{2}z^2+\frac{1}{6}z^3}{(1-z)^3}+\overline{\frac{\frac{1}{2}z^2+\frac{1}{6}z^3}{(1-z)^3}}.$$
This is the expected candidate for extremal mappings in the class $\mathcal{S}^0_\mathcal{H}$ of univalent harmonic mappings.

\subsection{Estimates of $\abs{a_2}$ of harmonic $K$-quasiconformal mappings}
We note that the problem of finding the sharp upper bound of $\abs{a_2}$ in the class $\mathcal{S}^0_\mathcal{H}(K)$
is also {\it challenging}. In 2022, Chen and Ponnusamy \cite[Proposition 1.6]{cp} proved the following rough estimate on $|a_2|$ for $f\in\mathcal{S}^0_\mathcal{H}(K)$.
\vskip.10in

\begin{thm} \label{thmA}
 For $K\geq1$, let $f=h+\overline{g}\in {\mathcal{S}_\mathcal{H}^0(K)}$, where $$h(z)=z +\sum_{n=2}^{\infty}a_n z^n\ {\textit and} \ g(z)=\sum_{n=2}^{\infty}b_n z^n.$$ Then 
$$|a_2|\leq\frac{32}{(1+1/K)^4}+\frac{64}{(1+1/K)^3}-2.$$
\end{thm}

\begin{rem}\label{r01}
{\rm For $K=1$ in Theorem \ref{thmA}, we observe that the result reduces to the case $\mathcal{S}^0_\mathcal{H}(1)=:\mathcal{S}$
of conformal mappings. 
At this time, we find that $$|a_2|\leq 8.$$ However, by Bieberbach's theorem, the sharp bound $$|a_2|\leq 2\ {\rm or}\ \alpha(\mathcal{S})=2$$ holds for $f\in\mathcal{S}$,  which shows that there still is substantial room for improvement on the upper bound of $\abs{a_2}$ for the family $\mathcal{S}^0_\mathcal{H}(K)$.}
\end{rem}





%\vskip.10in
%\noindent{\bf Theorem D.}
%\textit{If $f\in \mathcal{S}$, then $\|P_f\|\leq 6$ and $\|S_f\|\leq 6$.}

%\vskip.10in
%\noindent{\bf Theorem E.}
%\textit{If $\|P_f\|\leq 1$ and $\|S_f\|\leq 2$, then $f\in \mathcal{S}$.}


%\vskip.10in

%{\color{red}Can we generalize the above results of analytic case $\mathcal{S}$ into the case of harmonic $K$-quasiconformal case %$\mathcal{S}^0_\mathcal{H}(K)$?

%\vskip.10in
%\noindent{\bf Theorem F.}
%\textit{If $f\in \mathcal{S}^0_\mathcal{H}(K)$, then $\|P_f\|\leq \mu(K)$ and $\|S_f\|\leq \nu(K)$.}

%\vskip.10in
%\noindent{\bf Theorem G.}
%\textit{If $\|P_f\|\leq \phi(K)$ or $\|S_f\|\leq \psi(K)$, then $f\in \mathcal{S}^0_\mathcal{H}(K)$.}}

%\vskip.10in

%\subsection{The order of an affine and linear invariant family $\mathcal{F}_{\lambda}$}
%Let $\mathcal{F}_{\lambda}$ denote the subfamily of $\mathcal{F}$ which satisfies the condition 
%\begin{equation*}\label{sf}
%\|S_f\|\leq \lambda\quad (\lambda\geq0).
%\end{equation*}
%Also,
%we let $$\mathcal{F}_{\lambda}^0:=\{f\in\mathcal{F}_{\lambda}:\, g'(0)=0\}.$$

In 2017, Chuaqui, Hern\'{a}ndez, and Mart\'{i}n \cite[Theorem 1]{chm} proved the following result regarding the order
of the affine and linear invariant family $\mathcal{F}_{\lambda}$. 
We observe that this result will play an important role in our study.

\vskip.10in
\begin{thm} \label{thmD}
The order of the family $\mathcal{F}_{\lambda}$ is given by 
\begin{align*}
\begin{split}\alpha\left(\mathcal{F}_{\lambda}\right)=\sqrt{1+\frac{\lambda}{2}+\frac{1}{2}\sup\limits_{f\in\mathcal{F}_{\lambda}^0}
\abs{g''(0)}^2}+\frac{1}{2}\sup\limits_{f\in\mathcal{F}_{\lambda}^0}
\abs{g''(0)}. \end{split}
\end{align*}
\end{thm}


%\begin{theorem}\label{t11}
%Suppose that $f\in\mathcal{S}^0_\mathcal{H}(K)$. Then $$\|S_f\|\leq \psi(K).$$
%\end{theorem}


%\begin{theorem}\label{t1}
%\textit{Suppose that $f\in\mathcal{S}^0_\mathcal{H}(K)$}. Then $$\abs{a_2}\leq\frac{16}{\left(1+{1}/{K}\right)^2}-2.$$
%\end{theorem}

\vskip.20in
\section{Harmonic $K$-quasiconformal Koebe functions}\label{s3}

A long-standing bottleneck in the study of harmonic quasiconformal mappings lies in the absence of a canonical Koebe-type extremal function, which has significantly impeded the advancement of its extremal theory as well as the generalization of classical conformal mapping results to the harmonic quasiconformal framework. In this section, we remedy this critical gap by constructing such functions for the first time.

In what follows, we construct the harmonic $K$-quasiconformal Koebe functions
$$f_k(z)=h(z)+\overline{g(z)}$$
via the \textit{weighted shearing technique}, an approach adapted from the classical shearing method introduced by Clunie and Sheil-Small \cite{cs}.
The defining relations for its components are given by
$$h(z)-g(z)=\frac{z}{(1-z)^2}$$
and
$$\omega(z)=\frac{g'(z)}{h'(z)}=k z\quad (0\leq k<1).$$

%In this section, we construct the following so-called \textit{harmonic quasiconformal Koebe function} $$f_k(z)=h(z)+\overline{g(z)}$$ by using the {\it weight shearing method}, which originates from the classical shearing method proposed by Clunie and Sheil-Small \cite{cs}, where $$h(z)-g(z)=\frac{z}{(1-z)^2}$$ and $$\omega(z)=\frac{g'(z)}{h'(z)}=k\,z\quad (0\leq k<1).$$ 
%Theorem \ref{t3} given later will demonstrate its significance. %For more characterizations on other type of generalized harmonic Koebe functions, we refer the reader to the work of Ferrada-Salas and Mart\'{i}n \cite{fm}.



%\begin{example} \label{e1}
Solving the above system of differential equations and normalization, for $0\leq k<1$, we obtain \begin{align}\begin{split}\label{31}
  f_k(z)&=\frac{1}{(k-1)^3}{\left[\frac{(k-1)(1-3k+2kz)z }{(1-z)^2}+k(k+1) \log\left(\frac{1-z}{1-k z}\right)\right]}\\ & \qquad\ \ \ +\frac{k}{(k-1)^3}\overline{{\left[\frac{(1-k)(1+k-2z)z}{(1-z)^2}+(k+1) \log\left(\frac{1-z}{1-k z}\right)\right]} }.\end{split}
\end{align}
The power series form of $f_k$ is %presented as follows:
%  \begin{align*}f_k(z)&=z+\frac{1}{2}(k+4)z^2+\frac{1}{3}\left(k^2+4 k+9\right)z^3
%+\frac{1}{4}\left(k^3+4k^2+9k+16\right)z^4+\cdots
%\\&\quad\quad\ +\overline{\frac{k}{2}z^2+\frac{1}%{3}k(k+4)z^3+\frac{1}{4}k\left(k^2+4k+9\right)z^4
%+\cdots} \quad(0\leq k<1).
%\end{align*}
\begin{equation}\label{032}
\begin{aligned}
    f_k(z)=z+\sum_{n=2}^{\infty}A(n,k)z^n+\sum_{n=2}^{\infty}B(n,k)\overline{z}^n,
\end{aligned}
\end{equation}
where \begin{equation}\begin{aligned}\label{033}
A(n,k):=\frac{(1-k)^2n^2-2k(1-k)n+k(1+k)(1-k^n)}{(1-k)^3\,n}
\end{aligned}
\end{equation} and \begin{equation}\begin{aligned}\label{034}
B(n,k):=\frac{k(1-k)^2n^2-2k(1-k)n+k(1+k)(1-k^n)}{(1-k)^3\,n}.
\end{aligned}
\end{equation}

The images of the unit disk under analytic Koebe function $f_{0}(z)$, 
$f_{1/5}(z)$, $f_{2/5}(z)$, $f_{3/5}(z)$, $f_{4/5}(z)$ 
and harmonic Koebe function $\mathbb{K}(z)$ are illustrated, respectively, in Figure \ref{Fig_koebe_table}.

\begin{figure}[H] 
\centering 
\subfigure[Image of analytic Koebe function $f_0(z)$.]{ 
\label{Fig.sub.1} 
\includegraphics[height=6.0cm]{k1.jpg}} 
\subfigure[Image of $f_{1/5}(z)$.]{ 
\label{Fig.sub.2} 
\includegraphics[height=6.0cm]{k2.jpg}} 
%\label{Fig.lable} 
\end{figure}



\begin{figure}[H] 
\centering 
\subfigure[Image of $f_{2/5}(z)$.]{ 
\label{Fig.sub.3} 
\includegraphics[height=6.0cm]{k3.jpg}} 
\subfigure[Image of $f_{3/5}(z)$.]{ 
\label{Fig.sub.4} 
\includegraphics[height=6.0cm]{k4.jpg}} 
%\label{Fig.lable} 
\end{figure}


\begin{figure}[H] 
\centering 
\subfigure[Image of $f_{4/5}(z)$.]{ 
\label{Fig.sub.5} 
\includegraphics[height=6.0cm]{k5.jpg}} 
\subfigure[Image of harmonic Koebe function $\mathbb{K}(z)$.]{ 
\label{Fig.sub.6} 
\includegraphics[height=6.0cm]{k6.jpg}} 
\caption{\small{Images of the unit disk under the mappings $f_{0}(z)$, $f_{1/5}(z)$, $f_{2/5}(z)$, $f_{3/5}(z)$, $f_{4/5}(z)$, 
and $\mathbb{K}(z)$.}}
\label{Fig_koebe_table} 
\end{figure}




%The images of the concentric circle under $f_{1/2}(z)$ and $f_{2/3}(z)$ are illustrated in Figures \ref{fig3}
%and \ref{fig2}, respectively.
%\end{example}

%{\begin{figure}[ht]
%\centering
%\includegraphics[height=8.5cm]{12.pdf}
%\caption{Image of $f_{1/2}(z)$
%}\label{fig3}
%\end{figure}}

%{\begin{figure}[ht]
%\begin{center}
%\includegraphics[height=8.3cm]{13.pdf}
%\end{center}
%\caption{Image of $f_{2/3}(z)$
%}\label{fig2}
%\end{figure}}

By noting that the coefficient conjecture of Clunie and Sheil-Small \cite{cs} holds for the close-to-convex, starlike and typically real harmonic mappings, and harmonic quasiconformal mappings is a subclass of univalent harmonic mappings, the Bieberbach-type coefficient conjecture for \textbf{HQC} is confirmed valid by Li and Ponnusamy \cite{lp} for close-to-convex, starlike, and typically real harmonic quasiconformal mappings (see also \cite{dhr}). For recent investigations on Koebe-type functions in the setting of harmonic quasiconformal mappings, see \cite{wqr}.
%\end{rem}

%\begin{rem}\label{r03} {\rm 

It is natural to conjecture that the function $f_k(z)$ serves as the extremal mapping for the class $\mathcal{S}^0_\mathcal{H}(K)$ of \textbf{HQC} mappings. This not only opens a novel perspective on the class $\mathcal{S}^0_\mathcal{H}(K)$, but also provides a fundamental bridge connecting the theories of classical conformal mappings, harmonic mappings, quasiconformal mappings, Teichmüller spaces, and even SLE theory.

Indeed, the above formulas lead us to conjecture that the extremal function in $\mathcal{S}^0_\mathcal{H}(K)$ is unique and coincides with the harmonic $K$-quasiconformal Koebe functions.

\begin{con}\label{conj00}{\rm (\textbf{HQC Bieberbach conjecture})}
 For $K\geq1$, let $f=h+\overline{g}\in {\mathcal{S}_\mathcal{H}^0(K)}$, where $$h(z)=z +\sum_{n=2}^{\infty}a_n z^n\ {\rm and} \ g(z)=\sum_{n=2}^{\infty}b_n z^n.$$ Then 
$$\abs{a_n}\leq A(n,k),\, \abs{b_n}\leq B(n,k),\,{\rm and}\, \abs{\abs{a_n}-\abs{b_n}}\leq n\,\, {\rm for}\,\, n\geq 2,$$
where $A(n,k)$ and $B(n,k)$ are given by \eqref{033} and \eqref{034}, respectively, with $$k=\frac{K-1}{K+1}\quad(K\geq1).$$
\end{con}

However, this somewhat is premature in view that the related conjecture for the class $\mathcal{S}_\mathcal{H}^0$ (see e.g., \cite[pp. 87--88]{d}) remains 
an \textit{open} problem.
%\end{rem}

By the above coefficient conjecture for the class $\mathcal{S}^0_\mathcal{H}(K)$, we obtain the conjectured sharp upper bound of $\abs{a_2}$, which will play a distinguished role to study various characterizations for the class  $\mathcal{S}^0_\mathcal{H}(K)$ (or even the class $\mathcal{S}^0_\mathcal{H}$). 

\begin{con}\label{conj01}
Suppose that $f=h+\overline{g}\in\mathcal{S}^0_\mathcal{H}(K)$. Then
$$\sup\limits_{f\in\mathcal{S}_\mathcal{H}^0(K)}\abs{a_2}=\frac{5K+3}{2K+2}.$$
The bound is expected to be sharp.
\end{con}



By Conjecture \ref{conj01}, Nowak’s conjecture and Proposition \ref{prop1}, it is natural to state the following conjecture
associated with the sharp order of the family $\mathcal{S}_\mathcal{H}(K)$ to belong to a harmonic Hardy space $h^p$.

\begin{con}\label{conj1}
Suppose that $f=h+\overline{g}\in\mathcal{S}_\mathcal{H}(K)$. Then
$$\sup\limits_{f\in\mathcal{S}_\mathcal{H}(K)}\abs{a_2}=\frac{3K+1}{K+1},$$
and, in particular, %if $f$ preserves affine and linear invariant, 
we have %\begin{enumerate}
%\item for $\vartheta(K)\leq 2K$, we have 
$$f\in h^p\quad \left(0<p<\frac{1}{2K}\right).$$
%\vskip.05in
%\item for $\vartheta(K)>2K$, we have $$f\in h^p\quad \left(0<p<\frac{1}{\vartheta(K)}\right).$$
%\end{enumerate}
The bound for the order is expected to be sharp.
\end{con}


\begin{rem}\label{r0}
{\rm If one can prove Nowak’s
conjecture and the first part of Conjecture \ref{conj1}, it means that Problem \ref{p1} (Pavlovi\'{c}'s \textit{open problem}) can be \textit{solved} completely.}
\end{rem}





In what follows, we pose the conjecture associated with \textit{radius of covering theorem} of the family $\mathcal{S}^0_\mathcal{H}(K)$.
\begin{con} 
For $0<r<1$, $\D_r:=\{z: z\in\C\ {\it and}\ \abs z<r\}$ and  $K\geq1$, let $f\in {\mathcal{S}_\mathcal{H}^0(K)}$. Then 
$$R=\sup\limits_{0<r<1}\left\{r: \D_r\subseteq f(\D),\, f\in\mathcal{S}_\mathcal{H}^0(K)\right\}\geq\frac{K+1}{6K+2}.$$
The radius is expected to be sharp.
\end{con}

\begin{rem}\label{r3}
{\rm %For $k=1$, $f_1(z)$ reduces to the known \textit{harmonic Koebe function}; For $0\leq k<1$, we call $f_k(z)$ the \textit{harmonic quasiconformal Koebe function}. Also, 
By a direct calculation, we see that $$\|S_{f_k}\|\leq\frac{19}{2}
\quad(0\leq k<1).$$}
\end{rem}
%=============================
%$f(z)=h+\overline{g}$  where

%$h=\frac{\frac{(k-1) (k (z-2)+z)}{(z-1)^2}+k (k+1) \log (z-1)-k (k+1) \log (1-k z)}{(k-1)^3}$

%$g=\frac{k \left(-\frac{(k-1) z (k-2 z+1)}{(z-1)^2}+(k+1) \log (1-z)-(k+1) \log (1-k z)\right)}{(k-1)^3}$

%let $h,g$ be normorlized. we have 

%$h=z+ \frac{1}{2} (k+4) z^2+\frac{1}{3} \left(k^2+4 k+9\right) z^3+\frac{1}{4} \left(k^3+4 k^2+9 k+16\right) z^4+\frac{1}{5} \left(k^4+4 k^3+9 k^2+16 k+25\right) z^5+\frac{1}{6} \left(k^5+4 k^4+9 k^3+16 k^2+25 k+36\right) z^6+...$

%and 

%$$g=\frac{k z^2}{2} +\frac{1}{3} k (k+4) z^3+ \frac{1}{4} k \left(k^2+4 k+9\right) z^4+\frac{k \left(\frac{1}{5} (k+1) k^5+\frac{1}{5} (-k-1)+(k-1) (8-5 %(k+1))\right) z^5}{(k-1)^3}+\frac{k \left(\frac{1}{6} (k+1) k^6+\frac{1}{6} (-k-1)+(k-1) (10-6 (k+1))\right) z^6}{(k-1)^3}+$$

%==================================
We observe that Graf \cite[Theorem 3]{g}, Hern\'{a}ndez and Mart\'{i}n \cite[Theorem~5]{hm} once showed that $\|S_f\|$ are bounded for locally univalent and univalent harmonic mappings, respectively,
but the questions about the sharp bounds are still \textit{open}. By Conjecture \ref{conj01} and Remark \ref{r3}, it is natural to pose the following conjecture:

\begin{con}\label{cc1}
Suppose that $f=h+\overline{g}\in\mathcal{S}^0_\mathcal{H}(K)$. Then $$\|S_f\|\leq\frac{19K^2+26K+3}{2(K+1)^2}.$$
The bound is expected to be sharp.
\end{con}



\vskip.20in
\section{Proof of Theorem \ref{t1}}

%\subsection{Generalized Schwarz lemma}
We start by noting the following variant of the classical Schwarz lemma.
%For recent results involving Schwarz lemma or its generalizations, one can find in \cite{bkr,lt}. 

\begin{lemma} {\rm (\cite{b})}\label{t111}
Suppose that $0<k\leq1$ and $\omega$ is analytic in $\D$ with $$\omega(0)=0\ \textit{and}\ \abs{\omega(z)}\leq k.$$ Then
 $$\abs{\omega(z)}\leq k\abs{z}\ \textit{and}\ \abs{\omega'(0)}\leq k.$$ Equality holds for the functions $$\omega(z)=k\,e^{i\theta}z\quad(\theta\in[0,2\pi)).$$
\end{lemma}
%\begin{proof}
%By defining $$g(z):=\frac{\omega(z)}{z}\quad\left(z\in\D\setminus\{0\}\right)$$ and $$g(0)=\omega'(0),$$ we
%see that \begin{equation}\label{11}\abs{g(z)}\leq \frac{k}{r}\end{equation} on each circle $\abs z=r<1$. In view of the %maximum modulus
%theorem, we obtain that \eqref{11} also holds in $\abs {z}\leq r$. Let $r\rightarrow 1^{-}$ to get %$$\abs{\frac{\omega(z)}{z}}=\abs{g(z)}\leq
%k.$$ Subsequently, we have $\abs{\omega'(0)}\leq k$.
%\end{proof}

%\subsection{Estimating $\abs{b_2}$ of a harmonic $K$-quasiconformal mapping}
Now, we provide the sharp upper bound of $\abs{b_2}$ for the class $\mathcal{S}^0_\mathcal{H}(K)$, which will play a crucial role in determining the sharp upper bound of $\abs{a_2}$ of the family of harmonic $K$-quasiconformal mappings with bounded Schwarzian norm.

\begin{lemma}\label{t011}
Suppose that $f=h+\overline{g}\in\mathcal{S}^0_\mathcal{H}(K)$, where $$h(z)=z +\sum_{n=2}^{\infty}a_n z^n\ {\textit and} \ g(z)=\sum_{n=2}^{\infty}b_n z^n.$$ Then the sharp inequality \begin{equation}\label{22}\abs{b_2}\leq \frac{K-1}{2(K+1)}\end{equation}
holds.
\end{lemma}
\begin{proof}
For $f=h+\overline{g}\in\mathcal{S}^0_\mathcal{H}(K)$,
we see that its analytic dilatation $$\omega=\frac{g'}{h'}$$ satisfies the condition $$\abs{\omega(z)}\leq k\quad (0\leq k<1).$$
Moreover, if $f\in\mathcal{S}^0_\mathcal{H}(K)$, we know that $g'(0)=0$, i.e., $\omega(0)=0$. Thus, the assertion of
Lemma \ref{t111} shows that $$\abs{\omega'(0)}\leq k\quad (0\leq k<1).$$ By \cite[p. 87, Sect. 5.4]{d}, we get
 \begin{align*}
 \begin{split}\omega(z)&=\frac{g'(z)}{h'(z)}
 =\frac{2b_2z+3b_3z^2+\cdots}{1+2a_2z+\cdots}
=2b_2z+(3b_3-4a_2b_2)z^2+\cdots,
 \end{split}
\end{align*}
therefore, $$\abs{\omega'(0)}=\abs{2b_2}\leq k=\frac{K-1}{K+1},$$ which implies that the desired assertion of Lemma \ref{t011} is true. Equality holds for the functions $$\omega(z)=k\,e^{i\theta}z\quad\left(0\leq k<1;\, \theta\in[0,2\pi)\right).$$
This completes the proof of Lemma \ref{t011}.
\end{proof}
%\section{Proof of Theorem \ref{t2}}

In what follows, we introduce the \textit{strongly affine transformation} for mappings \( f = h + \overline{g} \in \mathcal{S}_{\mathcal{H}}(K) \), defined by
\[
\widetilde{A}_\xi(f)(z) := \frac{f(z) - \overline{\xi f(z)}}{1 - \overline{\xi} g'(0)},
\]
where the parameter \( \xi \) satisfies
\begin{equation}
\label{111}
|\xi| \leq \eta \quad \left(0\leq\eta< 1\right).
\end{equation}

 A family $\mathcal{\widetilde{F}}$ is said to be a \textit{strongly affine and linear invariant family}
(SAL family) if it is closed under both the Koebe transform and strongly affine transformation.

\begin{lemma}\label{t012}
The families \( \mathcal{S}_{\mathcal H}(K) \) and \({\mathcal{S}_\mathcal{H}}(K,\lambda)\) belong to $\mathcal{\widetilde{F}}$ of strongly affine and linear invariant family.
\end{lemma}
\begin{proof}
It easily to check that the family \( \mathcal{S}_{\mathcal H}(K) \) (resp. \({\mathcal{S}_\mathcal{H}}(K,\lambda)\)) is a linear invariant family (see \cite[p. 9]{DR}).
Let $\widetilde{A}_{\xi}(f)(z)=H(z)+\overline{G}(z)$, and denote its complex dilatation by $$\widetilde{\omega}(z)=\frac{G'(z)}{H'(z)}.$$

Note that %\(H(0)=H'(0)-1=G(0)=0\), and 
the constraint \[|\xi| \leq \frac{k - |\omega(z)|}{1 - k|\omega(z)|}<1 \quad \left(|\omega(z)| = \left|\frac{g'}{h'}\right| \leq k; \ 0\leq k < 1\right)\] on \(\xi\) is equivalent to
\[\left| \frac{\omega(z) - \xi}{1 - \overline{\xi} \omega(z)} \right| \leq k < 1.\] In conjunction with the definition of \(\widetilde{\omega}(z)\), we have \[|\widetilde{\omega}(z)| = \left| \frac{G'(z)}{H'(z)} \right| = \left| \frac{g' - \xi h'}{h' - \overline{\xi} g'} \right| = \left| \frac{\omega(z) - \xi}{1 - \overline{\xi} \omega(z)} \right| \leq k < 1,\]
which shows that the family \( \mathcal{S}_\mathcal{H}(K) \) (resp. \({\mathcal{S}_\mathcal{H}}(K,\lambda)\)) is a strongly affine invariant family. 
%Recall that 
%\(f \in \mathcal{S}_\mathcal{H}(K)\) (resp. \(f\in\widehat{\mathcal{S}_\mathcal{H}}(K,\lambda)\)) is a univalent harmonic quasiconformal mapping in \(\mathbb{D}\). %, and 
%by the classical property of a locally univalent quasiconformal mapping in \(\mathbb{D}\) with the modulus of complex dilatation less than 1 is univalent.
Thus, we deduce that the families \( \mathcal{S}_\mathcal{H}(K) \) and \({\mathcal{S}_\mathcal{H}}(K,\lambda)\) are {strongly affine and linear invariant families}. 
\end{proof}



We now turn our attention to the following result involving the order of the family \(\widehat{\mathcal{S}_\mathcal{H}}(K,\lambda)\).

\begin{prop} \label{c1}
Let $f\in\widehat{\mathcal{S}_\mathcal{H}}(K,\lambda)$. %such that $f$ preserves affine invariance. % preserve affine and linear invariant. 
Then
\begin{align*}
\sup\limits_{f\in\widehat{\mathcal{S}_\mathcal{H}}(K,\lambda)}\abs{a_2}=\sqrt{1+\frac{\lambda}{2}+\frac{1}{2}\left(\frac{K-1}{K+1}\right)^2}+\frac{K-1}{2\left({K+1}\right)}. 
\end{align*}
\end{prop}

\begin{proof} 
By harmonic quasiconformal Koebe function $f_k$ and Remark \ref{r3}, 
we see that $f_k\in \widehat{\mathcal{S}_\mathcal{H}}(K,\lambda)$, and  
inequality \eqref{22} also holds for this family. Note that Theorem \ref{thmD} remains valid for the strongly family \(\widehat{\mathcal{S}_\mathcal{H}}(K,\lambda)\).
By Lemmas \ref{t011} and \ref{t012}, and using the similar method as in Theorem \ref{thmD},
 we get the assertion of Proposition \ref{c1}.
\end{proof}


%Now, we establish the following relationship between
%the family $\widehat{\mathcal{S}^0_\mathcal{H}}(K,\lambda)$ of harmonic $K$-quasiconformal mappings with bounded %Schwarzian norm and harmonic Hardy spaces $h^p_K$.


%\begin{theorem}
%If $f\in\mathcal{S}^0_\mathcal{H}(K,\lambda)$ with $0\leq\lambda<6$, then
%$f\in h^p\, (0<p\leq 1/\varphi(K,\lambda))$,
%where $$\varphi(K,\lambda):=\alpha\left(\mathcal{S}^0_\mathcal{H}(K,\lambda)\right)
%=\sqrt{1+\frac{\lambda}{2}+\frac{1}{2}\left(\frac{K-1}{K+1}\right)^2}.$$
%\end{theorem}


In what follows, we are ready to give the proof of Theorem \ref{t1}. %which is presented in Sect. \ref{s1}.
\begin{proof}[Proof of Theorem \ref{t1}]
If $f\in\widehat{\mathcal{S}_\mathcal{H}}(K,\lambda)$, %preserve affine and linear invariant, 
it follows from Proposition \ref{c1} that
$$\sup\limits_{f\in\widehat{\mathcal{S}_\mathcal{H}}(K,\lambda)}\abs{a_2}=\sqrt{1+\frac{\lambda}{2}+\frac{1}{2}\left(\frac{K-1}{K+1}\right)^2}+\frac{K-1}{2\left({K+1}\right)}\geq1.$$
For convenience, we define a real function $\Phi(K)$ with one parameter $\lambda\, (\lambda\geq0)$ as follows: \begin{equation}\label{33}\Phi(K):=\sqrt{1+\frac{\lambda}{2}+\frac{1}{2}\left(\frac{K-1}{K+1}\right)^2}+\frac{K-1}{2\left({K+1}\right)}-2K\quad(K\geq1).\end{equation}

Now, we divide the proof into three cases.
\vskip.05in
(i) When $0\leq\lambda\leq6$ and $K\geq1$, we know that $\Phi(K)$ is a decreasing continuous function with respect to $K$
in $[1,+\infty)$,
and $$\Phi(1)=\sqrt{1+\frac{\lambda}{2}}-2\leq0,$$ which implies that $$\Phi(K)\leq0\quad(K\geq1).$$ Thus, we have
\begin{equation}\label{3}\varphi(K,\lambda)=\sqrt{1+\frac{\lambda}{2}+\frac{1}{2}\left(\frac{K-1}{K+1}\right)^2}+\frac{K-1}{2\left({K+1}\right)}\leq 2K,\end{equation}
that is, \begin{equation}\label{4}\frac{1}{\varphi(K,\lambda)}\geq\frac{1}{2K}.\end{equation}
By Proposition \ref{prop1},
it shows that $$f\in h^p\quad \left(0<p<\frac{1}{2K}\right).$$
\vskip.05in
(ii) When $\lambda>6$, we know that $$\Phi'(K)<0\quad (K\geq1).$$ Furthermore, we find that $$\Phi(1)=\sqrt{1+\frac{\lambda}{2}}-2>0$$
and $$\Phi\left(\lambda\right)<0.$$ Applying the zero point theorem to the real continuous function $\Phi(K)$ in bounded closed interval
$[1,\, \lambda]$, we deduce that the equation \begin{equation*}\sqrt{1+\frac{\lambda}{2}+\frac{1}{2}\left(\frac{K-1}{K+1}\right)^2}+\frac{K-1}{2\left({K+1}\right)}=2K\quad(K\geq1;\, \lambda>6),\end{equation*} 
which is equivalent to the equation \eqref{2}, exists a unique solution $K_1$ in $(1,+\infty)$.
If $K\geq K_1>1$, we know that
 \eqref{4} also holds. By Proposition \ref{prop1}, it implies that $$f\in h^p\quad \left(0<p<\frac{1}{2K}\right).$$
\vskip.05in
(iii) When $\lambda>6$ and $1\leq K<K_1$, we see that
\begin{equation*}\label{5}\varphi(K,\lambda)=\sqrt{1+\frac{\lambda}{2}+\frac{1}{2}\left(\frac{K-1}{K+1}\right)^2}+\frac{K-1}{2\left({K+1}\right)}>2K,\end{equation*}
or equivalently, $$\frac{1}{\varphi(K,\lambda)}<\frac{1}{2K}.$$ By Proposition \ref{prop1},
it means that $$f\in h^p\quad \left(0<p<\frac{1}{\varphi(K,\lambda)}\right).$$
The proof of Theorem \ref{t1} is thus completed.
\end{proof}



%\vskip.10in
%\section{Several examples}

%In this section, we will give three examples to illuminate Theorem \ref{t3}.



%\vskip.20in
%\section{Appendix}

%\noindent
%\textbf{Definition} When $\vert \xi \vert \leq \frac{k - \vert \alpha(\xi) \vert}{1 - k \vert \alpha(\xi) \vert}$ and $\vert \alpha(\xi) \vert \leq k < 1$, then $\frac{\vert \alpha(\xi) - \xi \vert}{\vert 1 - \overline{\alpha(\xi)} \xi \vert} \leq k$.


%\noindent
%\textbf{Lemma 1} When $\vert \xi \vert \leq \frac{k - \vert \alpha(\xi) \vert}{1 - k \vert \alpha(\xi) \vert}$ and $\vert \alpha(\xi) \vert \leq k < 1$, the function $f(t) = \frac{\vert \alpha(\xi) \vert + \vert \xi \vert^2 - \vert \alpha(\xi) \vert \vert \xi \vert^2}{1 + \vert \alpha(\xi) \vert \vert \xi \vert^2 - 2 \vert \alpha(\xi) \vert \vert \xi \vert t}$ is monotonically decreasing on $\mathbb{R}$.

%\medskip
%\noindent
%\textbf{Proof.} Since $0 < k < 1$, we have
%\[
%1 - k \vert \alpha(\xi) \vert - (k - \vert \alpha(\xi) \vert) = 1 - k - k \vert \alpha(\xi) \vert + \vert \alpha(\xi) \vert = 1 - k > 0.
%\]
%That is, $1 - k \vert \alpha(\xi) \vert > k - \vert \alpha(\xi) \vert \geq 0$, so $\frac{k - \vert \alpha(\xi) \vert}{1 - k \vert \alpha(\xi) \vert} < 1$, hence $\vert \xi \vert \leq \frac{k - \vert \alpha(\xi) \vert}{1 - k \vert \alpha(\xi) \vert} < 1$.

%Accordingly,
%\[
%f'(t) = \frac{2 \vert \alpha(\xi) \vert \vert \xi \vert \left( \vert \alpha(\xi) \vert + \vert \xi \vert^2 - \vert \alpha(\xi) \vert \vert \xi \vert^2 \right)}{\left( 1 + \vert \alpha(\xi) \vert \vert \xi \vert^2 - 2 \vert \alpha(\xi) \vert \vert \xi \vert t \right)^2} < 0.
%\]
%Thus, $f(t)$ is monotonically decreasing on $\mathbb{R}$.


%\noindent
%\textbf{Lemma 2} When $\vert \alpha(\xi) \vert < 1$, the function $g(u) = \frac{\vert \alpha(\xi) \vert + u}{1 + \vert \alpha(\xi) \vert u}$ is monotonically increasing on $\mathbb{R}$.

%\medskip
%\noindent
%\textbf{Proof.} Since $g'(u) = \frac{1 - \vert \alpha(\xi) \vert^2}{\left( 1 + \vert \alpha(\xi) \vert u \right)^2} > 0$, it follows that $g(u) = \frac{\vert \alpha(\xi) \vert + u}{1 + \vert \alpha(\xi) \vert u}$ is monotonically increasing on $\mathbb{R}$.


%\noindent
%\textbf{Proof} Let $\alpha(\xi) = \vert \alpha(\xi) \vert (\cos \alpha + i \sin \alpha)$, $\xi = \vert \xi \vert (\cos \theta + i \sin \theta)$. Then
%\[
%\frac{\vert \alpha(\xi) - \xi \vert}{\vert 1 - \overline{\alpha(\xi)} \xi \vert} = \frac{\left\vert \vert \alpha(\xi) \vert \cos \alpha - \vert \xi \vert \cos \theta + i \left( \vert \alpha(\xi) \vert \sin \alpha - \vert \xi \vert \sin \theta \right) \right\vert}{\left\vert 1 - \vert \alpha(\xi) \vert \vert \xi \vert \cos (\alpha - \theta) - i \vert \alpha(\xi) \vert \vert \xi \vert \sin (\alpha - \theta) \right\vert}
%\]

%\[= \sqrt{\frac{\left( \vert \alpha(\xi) \vert \cos \alpha - \vert \xi \vert \cos \theta \right)^2 + \left( \vert \alpha(\xi) \vert \sin \alpha - \vert \xi \vert \sin \theta \right)^2}{\sqrt{1 - 2 \vert \alpha(\xi) \vert \vert \xi \vert \cos (\alpha - \theta) + \left( \vert \alpha(\xi) \vert \vert \xi \vert \right)^2}}\]


%\[
%= \sqrt{\frac{\vert \alpha(\xi) \vert^2 + \vert \xi \vert^2 - 2 \vert \alpha(\xi) \vert \vert \xi \vert \cos (\alpha - \theta)}}{\sqrt{1 + \left( \vert \alpha(\xi) \vert \vert \xi \vert \right)^2 - 2 \vert \alpha(\xi) \vert \vert \xi \vert \cos (\alpha - \theta)}}.
%\]
%Since $\cos (\alpha - \theta) \in [-1, 1]$, by Lemma 1 we have
%\[
%f(\cos (\alpha - \theta)) \leq f(-1) = \frac{\left( \vert \alpha(\xi) \vert + \vert \xi \vert \right)^2}{1 + \left( \vert \alpha(\xi) \vert \vert \xi \vert \right)^2}.
%\]
%Because $\vert \xi \vert \leq \frac{k - \vert \alpha(\xi) \vert}{1 - k \vert \alpha(\xi) \vert}$, by Lemma 2 we get
%\[
%\frac{\vert \alpha(\xi) - \xi \vert}{\vert 1 - \overline{\alpha(\xi)} \xi \vert} \leq \sqrt{f(\cos (\alpha - \theta))} \leq \sqrt{\frac{\left( \vert \alpha(\xi) \vert + \vert \xi \vert \right)^2}{1 + \left( \vert \alpha(\xi) \vert \vert \xi \vert \right)^2}} = \frac{\vert \alpha(\xi) \vert + \vert \xi \vert}{1 + \vert \alpha(\xi) \vert \vert \xi \vert}
%\]
%\[
%= g(\vert \xi \vert) \leq g\left( \frac{k - \vert \alpha(\xi) \vert}{1 - k \vert \alpha(\xi) \vert} \right) = k.
%\]

%It should be noted that $\alpha(\xi) = k$ when $\vert \xi \vert = 0$.

%&=\frac{(16\alpha^2-24\alpha+11)\abs{\mu}^4+8(\alpha-1)(-4\alpha+1)\abs{\mu}^2+16\alpha(\alpha-1)}{2(1-\abs{\mu}^2)^2}.

%Let $x:=\abs{z}$ with $x\in(0,1)$. Now, we define
%\begin{equation*}
%\begin{aligned}
%p(x):=&\left[\frac{1}{\left(4-x^2\right)(1-x)}+\frac{3}{2}\frac{1}{\left(4-x^2\right)^2}\right]\left(1-x^2\right)^2\\
%=&\frac{2x^5+5x^4-10x^3-16x^2+8x+11}{2\left(4-x^2\right)^2}.
%\end{aligned}
%\end{equation*}
%By numerical computation (see Figure \ref{kx}), we confirm that
%\begin{equation}\label{416}
%\max\limits_{x\in(0,1)}\{p(x)\}\approx 0.38381.
%\end{equation}
%\begin{figure}[ht]
%\begin{center}
%\includegraphics[height=8.0cm]{kx.pdf}
%\end{center}
%\caption{\small{The graph of the function %$p(x)$.}}\label{kx}
%\end{figure}

\vskip .10in
	\noindent{\bf  Acknowledgements.}
%The authors thank the referees and editor for their valuable comments and suggestions,
%which were essential in improving the quality of this paper. 
The authors %also 
thank %Shaolin Chen,
Daoud Bshouty, 
%Xiaokai He, 
David Kalaj, %Shengyu Li, 
Gang Liu,
Zhi-Hong Liu, Saminathan Ponnusamy, 
Toshiyuki Sugawa,  %Miaokun Wang,
Yu-Dong Wu, %Gengrong Zhang, 
and 
 Jian-Feng Zhu
for their helpful conversations and useful comments in different stages of the preparation of this paper.


\vskip .10in
	\noindent{\bf  Funding.}
Z.-G. Wang was partially supported by the \textit{Key Project of Education Department of Hunan Province} under Grant no. 25A0668, and
the \textit{Natural Science Foundation of Changsha} under Grant no. kq2502003
of the P. R. China. X.-Y. Wang was partially supported by the \textit{National Scholarship Council of China} under Grant
no. 202306840137 of the P. R. China.
A. Rasila was partially supported by \textit{Natural Science Foundation of Guangdong Province} under Grant no. 2024A1515010467 of the P. R. China, and \textit{Li Ka Shing Foundation} under Grant no. 2024LKSFG06. 
J.-L. Qiu was partially supported by the \textit{Graduate Research Innovation Project of Hunan Province} under Grant no. CX20240803 of the P. R. China.


%The author would like to thank Prof. Antti Rasila and Dr. Zhi-Hong Liu for their fruitful
%and enlightening discussions in the preparation of this paper.
%Z.-H. Liu was supported by
%the \textit{National
%Natural Science Foundation} under Grant no. 11961013 of the P. R. China.

\vskip .10in
\noindent{\bf Conflicts of interest.} The authors declare that they have no conflict of interest.


\vskip .10in
\noindent{\bf Data availability statement.}  Data sharing is not applicable to this article as no datasets were generated or analysed during the current study.
\vskip.05in
\begin{thebibliography}{99}
\bibitem{aap}
{\small Y. Abu-Muhanna, R. M. Ali, and S. Ponnusamy, The spherical metric and univalent harmonic mappings,
\textit{Monatsh. Math.} \textbf{188} (2019), 703--716.}

%\vskip.05in
%\bibitem{Ahlfors-Weill-1962}
%{\small L. Ahlfors and   G. Weill,   A uniqueness theorem for Beltrami equation,  
% {\it Proc. Amer. Math. Soc.}  {\bf 13}  (1962),  975--978.}



%\vskip.05in
%\bibitem{Aghalary-Orouji-2014}
%R. Aghalary and  Z. Orouji, Norm estimates of the pre-Schwarzian derivatives for $\alpha$-spiral-like functions of order $\rho$. 
% {\it  Complex Anal. Oper. Theory}   {\bf 8} (2014), no. 4, 791--801.


%\vskip.05in
%\bibitem{acs}
%{\small H. Arbel\'{a}ez, M. Chuaqui, and W. Sierra, Nehari-type families of harmonic mappings, \textit{Math. Nachr.} \textbf{293} (2020), 39--51.}

%\vskip.05in
%\bibitem{ahs}
% {\small H. Arbel\'{a}ez, R. Hern\'{a}ndez, and W. Sierra, Lower and upper order of harmonic mappings, \textit{J. Math. Anal. Appl.} \textbf{507} (2022), Paper No. 125837, 18 pp.}

\vskip.05in
\bibitem{al}
{\small Y. Abu-Muhanna and A. Lyzzaik, The boundary behavior of harmonic univalent maps, \textit{Pacific J. Math.} \textbf{141} (1990), 1--20.}


\vskip.05in
\bibitem{am}
{\small A. Aleman and M. J. Mart\'{\i}n, Convex harmonic mappings are not necessarily in $h^{1/2}$, \textit{Proc. Amer. Math. Soc.} \textbf{143} (2015), 755--763.}




%\vskip.05in
%\bibitem{ap}
%{\small M. F. Ali and S. Pal, The Schwarzian norm estimates for Janowski convex functions, \textit{Proc. Edinburgh Math. Soc.} (2024). doi: 10.1017/S0013091524000014.}




%\vskip.05in
%\bibitem{ahs}
%{\small H. Arbel\'{a}ez, R. Hern\'{a}ndez and W. Sierra, Lower and upper order of harmonic mappings, \textit{J. Math. Anal. Appl.} \textbf{507} (2022), Paper No. 125837, 18 pp.}

\vskip.05in
\bibitem{ak}
{\small K. Astala and P. Koskela, $H^p$-theory for quasiconformal mappings, \textit{Pure Appl. Math. Q.} \textbf{7} (2011), 19--50.}





%\vskip.05in
%\bibitem{bkl}
%{\small S. Benedict, P. Koskela, and X. Li, Weighted Hardy spaces of quasiconformal mappings, \textit{J. Geom. Anal.} \textbf{32} (2022), Paper No. 97, 23 pp.}


%\vskip.05in
%\bibitem{Becker-1972}
%J. Becker, L\"ownersche Differentialgleichung und quasikonform fortsetzbare schlichte Funktionen, 
%\textit{J. Reine Angew. Math.} \textbf{255} (1972), 23--43.

\vskip.05in
\bibitem{b}
{\small M. Borovikov, On Koebe radius and coefficients estimate for univalent harmonic mappings, 
\textit{J. Anal.} \textbf{33} (2025), 2275--2284.}


%\vskip.05in
%\bibitem{bk}
%{\small O. Bouchala and P. Koskela, Existence of quasiconformal mappings in a given Hardy space, \textit{Proc. Amer. Math. Soc.}
%\textbf{152} (2024), 177--191.}


%\vskip.05in
%\bibitem{bkr}
%{\small F. Bracci, D. Kraus, and O. Roth, A new Schwarz-Pick lemma at the boundary and rigidity of holomorphic maps,
%\textit{Adv. Math.} \textbf{432} (2023), Paper No. 109262, 41 pp.}

%\vskip.05in
%\bibitem{bp}
%{\small S. V. Bharanedhar and S. Ponnusamy, Coefficient conditions for harmonic
%univalent mappings and hypergeometric mappings, \textit{Rocky Mountain J. Math.} \textbf{44}
%(2014), 753--777.}


%\vskip.05in
%\bibitem{cf}
%X. Chen and A. Fang, A Schwarz-Pick inequality for harmonic quasiconformal mappings and its applications, \textit{J. Math. Anal. Appl.} 
%\textbf{369} (2010), 22--28.

%\vskip.05in
%\bibitem{cp2}
%{\small S. Chen and S. Ponnusamy, John disks and $K$-quasiconformal harmonic mappings, \textit{J. Geom. Anal.} \textbf{27} (2017), 1468--1488.}

%\vskip.05in
%\bibitem{cp3}
%{\small S. Chen and S. Ponnusamy, Radial length, radial John disks and $K$-quasiconformal harmonic mappings, \textit{Potential Anal.} \textbf{50} (2019), 415--437.}

\vskip.05in
\bibitem{cp}
{\small S. Chen and S. Ponnusamy, Koebe type theorems and pre-Schwarzian of harmonic $K$-quasiconformal mappings, and their applications, \textit{Acta Math. Sin. (Engl. Ser.)} \textbf{38} (2022), 1965--1980.}

%\vskip.05in
%\bibitem{cpw}
%{\small S. Chen, S. Ponnusamy and X. Wang, Integral means and coefficient estimates on planar harmonic mappings, \textit{Ann. Acad. Sci. Fenn. Math.} \textbf{37} (2012), 69--79.}

\vskip.05in
\bibitem{cdo}
{\small M. Chuaqui, P. Duren, and B. Osgood, The Schwarzian derivative for harmonic mappings, \textit{J. Anal. Math.} \textbf{91} (2003), 329--351.}


\vskip.05in
\bibitem{chm}
{\small M. Chuaqui, R. Hern\'{a}ndez, and M. J. Mart\'{i}n, Affine and linear invariant families of harmonic mappings, \textit{Math. Ann.} \textbf{367} (2017), 1099--1122.}

\vskip.05in
\bibitem{cs}
{\small J. G. Clunie and T. Sheil-Small, Harmonic univalent functions, \textit{Ann. Acad. Sci. Fenn. Ser. A. I Math.} \textbf{9} (1984), 3--25.}


\vskip.05in
\bibitem{ds}
{\small S. Das and A. Sairam Kaliraj, Integral mean estimates for univalent and locally univalent harmonic mappings,
\textit{Canad. Math. Bull.} \textbf{67} (2024), 655--669.}


\vskip.05in
\bibitem{de}
{\small L. de Branges, A proof of the Bieberbach conjecture, \textit{Acta Math.} \textbf{154} (1985), 137--152.}

%\vskip.05in
%\bibitem{du70}
%{\small P. L. Duren, \textit{Theory of $H^p$ spaces}, Pure and Applied Mathematics, 38, Academic Press, New York,
%1970.}

\vskip.05in
\bibitem{dhr}
{\small S. Das, J. Huang, and A. Rasila, Hardy spaces of harmonic quasiconformal mappings and Baernstein's theorem, \textit{Bull. Sci. Math.} \textbf{208} (2026), 103789.}

\vskip.05in
\bibitem{DR} 
{\small S. Das, A.Rasila,
On harmonic quasiregular mappings in Bergman spaces,
\textit{Potential Anal.} \textbf{64} (2026), 26.}


\vskip.05in
\bibitem{du}
{\small P. L. Duren, \textit{Univalent functions}, Grundlehren der Mathematischen Wissenschaften, Vol. 259.
Springer, New York, 1983.}

\vskip.05in
\bibitem{d}
{\small P. Duren, \textit{Harmonic mappings in the plane}, Cambridge University
Press, Cambridge, 2004.}


%\vskip.05in
%\bibitem{eay}
%{\small A. Ebadian, S. Azizi and S. Yal\c{c}in, Univalent harmonic mappings and Hardy spaces, \textit{Turkish J. Math.} \textbf{43} (2019), 284--292.}

%\vskip.05in
%\bibitem{e}
%{\small
%I. Efraimidis, Criteria for univalence and quasiconformal extension for harmonic mappings on planar domains, \textit{Ann. Fenn. Math.}
%\textbf{46} (2021), 1123--1134.}

%\vskip.05in
%\bibitem{ehm}
%{\small I. Efraimidis, R. Hern\'{a}ndez and M. J. Mart\'{i}n, Ahlfors-Weill extensions for harmonic mappings, \textit{J. Math. Anal. Appl.} \textbf{523} (2023), Paper No. 127053, 15 pp.}

%\vskip.05in
%\bibitem{fh}
%{\small D. Fu and X. Huang, Harmonic $K$-quasiconformal mappings from unit disk onto half planes, \textit{Bull. Malays. Math. Sci. Soc.} \textbf{39} (2016), 339--347.}

\vskip.05in
\bibitem{g}
{\small S. Y. Graf, On the Schwarzian norm of harmonic mappings, \textit{Probl. Anal. Issues Anal.} \textbf{5(23)} (2016), 20--32.}


%\vskip.05in
%\bibitem{fm}
%{\small \'{A}. Ferrada-Salas and M. J. Mart\'{i}n,  Generalized harmonic Koebe functions, 
%\textit{J. Math. Anal. Appl.} \textbf{435} (2016), 860--873.}

%\vskip.05in
%\bibitem{Fait-Krzyz-Zygmunt-1976}
%{\small M. Fait,  J. G. Krzy\.{z}, and J.  Zygmunt, Explicit quasiconformal extensions for some classes of univalent functions,
% {\it  Comment. Math. Helv.}  {\bf  51} (1976), 279--285.}

%\vskip.05in
%\bibitem{gk}
%{\small A. Gjokaj and D. Kalaj, Quasiconformal harmonic mappings between the unit ball and a spatial domain with $C^{1,\alpha}$ boundary, 
%\textit{Potential Anal.} \textbf{57} (2022), 367--377.}

%\vskip.05in
%\bibitem{hss}
%J. Hubbard, D. Schleicher and S. Sutherland, How to find all roots of complex polynomials by Newton's method, \textit{Invent. Math.} \textbf{146} %(2001), 1--33.

%\vskip.05in
%\bibitem{hm0}
%R. Hern\'{a}ndez and M. J. Mart\'{i}n, Stable geometric properties of analytic and harmonic functions, 
%\textit{Math. Proc. Cambridge Philos. Soc.} \textbf{155} (2013), 343--359.

%\vskip.05in
%\bibitem{hk}
%{\small J. Ha and D. Kim, Maximal graphs and harmonic mappings, \textit{J. Math. Anal. Appl.} \textbf{527} (2023), Paper No. 127412, 17 pp.}

\vskip.05in
\bibitem{hm}
{\small R. Hern\'{a}ndez and M. J. Mart\'{i}n, Pre-Schwarzian and Schwarzian derivatives of harmonic mappings,
\textit{J. Geom. Anal.} \textbf{25} (2015), 64--91.}



%\vskip.05in
%\bibitem{j}
%{\small
%J. Jin, On a multiplier operator induced by the Schwarzian derivative of univalent functions,
%\textit{Bull. London Math. Soc.}  (2024), doi: 10.1112/blms.13056.}


%\vskip.05in
%\bibitem{hm22}
%{\small R. Hern\'{a}ndez and M. J. Mart\'{i}n, On the harmonic M\"{o}bius transformations, \textit{J. Geom. Anal.} \textbf{32} (2022), Paper No. 18, 27 pp.}

%\vskip.05in
%\bibitem{hf}
%{\small Z. Hu and J. Fan, Criteria for univalency and quasiconformal extension for harmonic mappings,
%\textit{Kodai Math. J.} \textbf{44} (2021), 273--289.}



%\vskip.05in
%\bibitem{h}
%{\small X. Huang, Harmonic quasiconformal mappings on the upper half-plane, \textit{Complex Var. Elliptic Equ.} \textbf{58} (2013), 1005--1011.}


\vskip.05in
\bibitem{ka}
{\small D. Kalaj, Muckenhoupt weights and Lindel\"{o}f theorem for harmonic mappings, \textit{Adv. Math.}
\textbf{280} (2015), 301--321.}

\vskip.05in
\bibitem{ka2}
{\small D. Kalaj, On Riesz type inequalities for harmonic mappings on the unit disk, \textit{Trans. Amer. Math. Soc.} \textbf{372} (2019), 4031--4051.}


%\vskip.05in
%\bibitem{ka2}
%{\small D. Kalaj, On Riesz type inequalities for harmonic mappings on the unit disk, \textit{Trans. Amer. Math. Soc.} \textbf{372} (2019), 4031--4051.}

%\vskip.05in
%\bibitem{Kanas-Sugawa-2011}
%{\small S. Kanas  and T. Sugawa,  Sharp norm estimate of Schwarzian derivative for a class of convex functions, {\it Ann. Polon. Math.} {\bf 101} (2011), 75--86.}
 
%\vskip.05in
%\bibitem{kmp}
%{\small S. Kanas, S. Maharana and J. K. Prajapat, Norm of the pre-Schwarzian derivative, Bloch's constant and coefficient bounds in some classes of harmonic mappings, \textit{J. Math. Anal. Appl.} \textbf{474} (2019), 931--943.}

%\vskip.05in
%\bibitem{kpv}
%D. Kalaj, S. Ponnusamy and M. Vuorinen, Radius of close-to-convexity and fully starlikeness of harmonic mappings,
%\textit{Complex Var. Elliptic Equ.} \textbf{59} (2014), 539--552.


%\vskip.05in
%\bibitem{kn}
%D. Khavinson and G. Neumann, From the fundamental theorem of algebra to astrophysics: a ``harmonious" path, \textit{Notices Amer. Math. Soc.} \textbf{55} (2008), 666--675.

%\vskip.05in
%\bibitem{Kuhnau-1971}
%R. K\"{u}hnau, Verzerrungss\"{a}tze und Koeffizientenbedingungen vom Grunskyschen Typ f\"{u}r quasikonforme Abbildungen. 
% (in German) {\it  Math. Nachr.}  {\bf 48} (1971), 77--105.

%\vskip.05in
%\bibitem{Kraus-1932}
%W. Kraus, {\"U}ber den Zusammenhang einiger Charakteristiken eines einfach zusammenh{\"a}ngenden Bereiches mit der %Kreisabbildung. 
% (in German) {\it  Mitt. Math. Sem. Giessen}  {\bf 21} (1932), 1--28.

%\vskip.05in
%\bibitem{lp}
%{\small G. Liu and S. Ponnusamy, Uniformly locally univalent harmonic mappings associated with the pre-Schwarzian norm, \textit{Indag. Math. (N.S.)} \textbf{29} (2018), 752--778.}

%\vskip.05in
%\bibitem{lt}
%{\small T. Liu and X. Tang, Schwarz lemma at the boundary of strongly pseudoconvex domain in $\Bbb{C}^n$, \textit{Math. %Ann.} \textbf{366} (2016), 655--666.}



\vskip.05in
\bibitem{ln}
{\small
J. Laitila, P. J. Nieminen, E. Saksman, and  H.-O. Tylli, Rigidity of composition operators on the Hardy space $H^p$, \textit{Adv. Math.} \textbf{319} (2017), 610--629.}

\vskip.05in
\bibitem{lp}
{\small
P. Li and S. Ponnusamy, On the coefficients estimate of $K$-quasiconformal harmonic mappings, 
\textit{Res. Math. Sci.} \textbf{13} (2026), Paper No. 16. }


\vskip.05in
\bibitem{lz}
{\small J. Liu and J.-F. Zhu, Riesz conjugate functions theorem for harmonic quasiconformal mappings, \textit{Adv. Math.} \textbf{434} (2023), Paper No. 109321, 27 pp.}

%\vskip.05in
%\bibitem{lp2}
%{\small Z.-H. Liu and S. Ponnusamy, Bohr radius for subordination and $K$-quasiconformal harmonic mappings, \textit{Bull. Malays. Math. Sci. Soc.}
%\textbf{42} (2019), 2151--2168.}

%\vskip.05in
%\bibitem{lw}
%{\small L. Liao and C. Wu, Exact meromorphic solutions of Schwarzian differential equations, \textit{Math.
%Z.} \textbf{300} (2022), 1657--1672.}



%\vskip.05in
%\bibitem{m}
%{\small O. Martio, On harmonic quasiconformal mappings, \textit{Ann. Acad. Sci. Fenn. Ser. A I}
%\textbf{425} (1968), 10 pp.}


%\vskip.05in
%\bibitem{ny}
%L. Nie and Z. Yang, The Schwarzian derivative of harmonic mappings in the plane, \textit{Chinese Ann. Math. Ser. B} \textbf{41} (2020), 193--208.

%\vskip.05in
%\bibitem{Nehari-1949}
%Z. Nehari, The Schwarzian derivative and schlicht functions, 
% {\it Bull. Amer. Math. Soc.} {\bf 55} (1949), 545--551.

\vskip.05in
\bibitem{mb}
{\small P. Melentijevi\'{c} and V. Bo\v{z}in,  Sharp Riesz-Fej\'{e}r inequality for harmonic Hardy spaces, \textit{Potential Anal.} \textbf{54} (2021), 575--580.}

\vskip.05in
\bibitem{n}
{\small M. Nowak, Integral means of univalent harmonic maps, \textit{Ann. Univ. Mariae Curie-Skłodowska Sect. A} \textbf{50} (1996), 155--162.}

\vskip.05in
\bibitem{psz}
{\small D. Partyka, K. Sakan, and J.-F. Zhu, Quasiconformal harmonic mappings with the convex holomorphic part, \textit{Ann. Acad. Sci. Fenn. Math.}
\textbf{43} (2018), 401--418.}

\vskip.05in
\bibitem{pav}
{\small M. Pavlovi\'{c}, \textit{Function classes on the unit disc: an introduction}, De Gruyter Studies in Mathematics, Vol. 52, De Gruyter, Berlin, 2014.}


%\vskip.05in
%\bibitem{prv}
%{\small F. P\'{e}rez-Gonz\'{a}lez, J. R\"{a}tty\"{a} and T. Vesikko, Integral means of derivatives of univalent functions in Hardy
%spaces, \textit{Proc. Amer. Math. Soc.} \textbf{151} (2023), 611--621.}


\vskip.05in
\bibitem{p}
{\small C. Pommerenke, Linear-invariante Familien analytischer Funktionen I, \textit{Math. Ann.} \textbf{155} (1964), 108--154.}


\vskip.05in
\bibitem{pqw}
{\small S. Ponnusamy, J. Qiao, and X. Wang, Uniformly locally univalent harmonic mappings, \textit{Proc. Indian Acad. Sci. Math. Sci.} \textbf{128} (2018), Paper No. 32, 14 pp.}


%\vskip.05in
%\bibitem{sy}
%{\small R. Schoen and S.-T. Yau, On univalent harmonic maps between surfaces, \textit{Invent. Math.} \textbf{44} (1978), 265--278.}



%\vskip.05in
%\bibitem{sz}
%{\small O. S\`{e}te and J. Zur, A Newton method for harmonic mappings in the plane, \textit{IMA J. Numer. Anal.} \textbf{40} (2020), 2777--2801.}


\vskip.05in
\bibitem{s}
{\small T. Sheil-Small, Constants for planar harmonic mappings, \textit{J. London Math.
Soc.} \textbf{42} (1990), 237--248.}

%\vskip.05in
%\bibitem{Sugawa-1997}
%T. Sugawa, On the norm of pre-Schwarzian derivatives of strongly starlike functions, %{\it  {S\={u}rikaisekikenky\=%{u}sho K\={o}ky\={u}roku}}   {\bf  1012} (1997), 178--185. 
% {\it Ann. Univ. Mariae Curie-Sk{\l}odowska, Sect. A} {\bf  52} (1998), 149--157.

%\vskip.05in
%\bibitem{Suita-1996}
% N. Suita,  Schwarzian derivatives of convex functions. 
 % {\it  J. Hokkaido Univ. Ed. Sect. II A} {\bf 46} (1996), 113--117.




\vskip.05in
\bibitem{srj}
{\small Y. Sun, A. Rasila, and Y.-P. Jiang, Linear combinations of harmonic quasiconformal mappings convex
in one direction, \textit{Kodai Math. J.}  \textbf{39} (2016), 366--377.}


\vskip.05in
\bibitem{t}
{\small V. Todor\v{c}evi\'{c}, \textit{Harmonic quasiconformal mappings and hyperbolic type metrics}, Springer, Cham, 2019.}


%\vskip.05in
%\bibitem{wlf}
%{\small X. Wang, H. Li and J. Fan, Pre-Schwarzian and Schwarzian norm estimates for subclasses of univalent functions,  
%\textit{Monatsh. Math.} (2024). doi: %10.1007/s00605-024-01971-1.}



%\vskip.05in
%\bibitem{wlrs}
%{\small Z.-G. Wang, Z.-H. Liu, A. Rasila, and Y. Sun, On a problem of Bharanedhar and Ponnusamy involving planar harmonic mappings,
%\textit{Rocky Mountain J. Math.} \textbf{48} (2018), 1345--1358.}

\vskip.05in
\bibitem{wqr}
{\small Z.-G. Wang, J.-L. Qiu, and A. Rasila, On Koebe-type functions for harmonic quasiconformal mappings, to appear in \textit{Anal. Math.} 2026.}

\vskip.05in
\bibitem{wsj}
{\small Z.-G. Wang, L. Shi, and Y.-P. Jiang, On harmonic $K$-quasiconformal mappings associated with asymmetric vertical strips,
\textit{Acta Math. Sin. (Engl. Ser.)} \textbf{31} (2015), 1970--1976.}


%\vskip.05in
%\bibitem{Yamashita-1999}
%S. Yamashita, Norm estimates for function starlike or convex of order alpha. 
% {\it  Hokkaido Math. J.}  {\bf  28} (1999), 217--230.


%\vskip.05in
%\bibitem{wlrs}
%Z.-G. Wang, Z.-H. Liu, A. Rasila and Y. Sun, On a problem of Bharanedhar and Ponnusamy involving planar harmonic mappings, \textit{Rocky Mountain J. Math.} {\bf 48} (2018), 1345--1358.

%\vskip.05in
%\bibitem{z}
%{\small J.-F. Zhu, Coefficients estimate for harmonic $v$-Bloch mappings and harmonic $K$-quasiconformal mappings, \textit{Bull. Malays. Math. Sci. Soc.}
% \textbf{39} (2016), 349--358.}


%\vskip.05in
%\bibitem{wsj}
%Z.-G. Wang, L. Shi and Y.-P. Jiang, On harmonic $K$-quasiconformal mappings
%associated with asymmetric vertical strips, \textit{Acta Math. Sin.} (\textit{Engl. Ser.}) \textbf{31} (2015), 1970--1976.


\end{thebibliography}

\end{document}

%------------------------------------------------------------------------------
% End of journal.tex
%------------------------------------------------------------------------------
